%%%%%%%%%%%%%%%%%%%%%%%%%%%%%%%%%%%%%%%%%%%%%%%%%%%
\begin{frame}[fragile]\frametitle{Convolution by Hand}
\begin{itemize}
\item Slide the 3 x 3 filter over the 5 x 5 image; at each position multiply matching entries and add
\item This filter responds where brightness jumps from left to right
\item Top-left position: $0\cdot(-1) + 0\cdot 0 + 1\cdot 1 = 1$ in each of 3 rows, total 3
\item Output size: $(N - F + 2P)/S + 1$ for image width $N$, filter width $F$, padding $P$, stride $S$
\item Here $N=5$ and $F=3$: stride 1 gives 3 x 3, stride 2 gives 2 x 2, padding 1 keeps 5 x 5
\end{itemize}
\begin{center}
\footnotesize
$\begin{bmatrix}0&0&1&1&1\\0&0&1&1&1\\0&0&1&1&1\\0&0&1&1&1\\0&0&1&1&1\end{bmatrix}
\ast
\begin{bmatrix}-1&0&1\\-1&0&1\\-1&0&1\end{bmatrix}
=
\begin{bmatrix}3&3&0\\3&3&0\\3&3&0\end{bmatrix}$
\end{center}
\end{frame}

%%%%%%%%%%%%%%%%%%%%%%%%%%%%%%%%%%%%%%%%%%%%%%%%%%
\begin{frame}[fragile]\frametitle{Quick Check: CNN Concepts}
\begin{block}{Think About It}
1. A 28 x 28 image goes through one 3 x 3 filter, stride 1, no padding. What size is the feature map?
\begin{enumerate}
  \item[A.] 28 x 28
  \item[B.] 26 x 26
  \item[C.] 25 x 25
  \item[D.] 9 x 9
\end{enumerate}
2. Why does a CNN need far fewer parameters than a fully connected network on the same image?
\begin{enumerate}
  \item[A.] It ignores most of the pixels.
  \item[B.] The same small filter is reused at every position.
  \item[C.] Pooling deletes pixels before training.
  \item[D.] It trains with a smaller learning rate.
\end{enumerate}
\end{block}
\end{frame}

%%%%%%%%%%%%%%%%%%%%%%%%%%%%%%%%%%%%%%%%%%%%%%%%%%
\begin{frame}[fragile]\frametitle{Quick Check: CNN Concepts}
\begin{block}{Answer}
\textbf{1. Correct: B.} $(28 - 3)/1 + 1 = 26$. A would need padding of 1, C forgets the $+1$, and D divides by the filter size.

\textbf{2. Correct: B.} Weight sharing is the main saving. Every pixel is still seen by some window, so A is wrong; pooling shrinks later layers but does not create the saving, and the learning rate has no effect on the parameter count.
\end{block}
\end{frame}

%%%%%%%%%%%%%%%%%%%%%%%%%%%%%%%%%%%%%%%%%%%%%%%%%%
\begin{frame}[fragile]\frametitle{}
\begin{center}
{\Large Case Study: Reading Handwritten Digits}
\end{center}
\end{frame}

%%%%%%%%%%%%%%%%%%%%%%%%%%%%%%%%%%%%%%%%%%%%%%%%%%%%%%%%%%
\begin{frame}[fragile]\frametitle{The Data}
\begin{itemize}
\item MNIST: 60,000 training and 10,000 test images of handwritten digits, 28 x 28 pixels, grey levels scaled to 0 to 1
\item The same task shows up in engineering work: reading meter and gauge displays, part serial numbers, batch codes stamped on components
\item Ten classes, roughly balanced: from 5,421 images of the digit 5 to 6,742 of the digit 1
\item PyTorch wants images as (N, channels, height, width)
\end{itemize}
\begin{lstlisting}
from torchvision import datasets

train = datasets.MNIST('data', train=True, download=True)
xtr = train.data.float().div(255).unsqueeze(1)

>>> tuple(xtr.shape)
(60000, 1, 28, 28)
\end{lstlisting}
\end{frame}

%%%%%%%%%%%%%%%%%%%%%%%%%%%%%%%%%%%%%%%%%%%%%%%%%%%%%%%%%%
\begin{frame}[fragile]\frametitle{A Small CNN}
\begin{itemize}
\item Two convolution layers, 8 and 16 filters of 3 x 3, each followed by ReLU and 2 x 2 max pooling
\item One fully connected layer turns the 400 numbers left into 10 scores, one per digit
\item The softmax sits inside the loss function, so the network outputs raw scores
\end{itemize}
\begin{lstlisting}
class SmallCNN(nn.Module):
    def __init__(self):
        super().__init__()
        self.conv1 = nn.Conv2d(1, 8, kernel_size=3)
        self.conv2 = nn.Conv2d(8, 16, kernel_size=3)
        self.fc = nn.Linear(16 * 5 * 5, 10)

    def forward(self, x):
        x = F.max_pool2d(F.relu(self.conv1(x)), 2)
        x = F.max_pool2d(F.relu(self.conv2(x)), 2)
        return self.fc(x.flatten(1))
\end{lstlisting}
\end{frame}

%%%%%%%%%%%%%%%%%%%%%%%%%%%%%%%%%%%%%%%%%%%%%%%%%%%%%%%%%%
\begin{frame}[fragile]\frametitle{Shapes Through the Network}
\begin{itemize}
\item A 3 x 3 filter with no padding trims one pixel from each edge: 28 becomes 26
\item Pooling halves the size and rounds down: 11 becomes 5
\item 16 x 5 x 5 = 400 numbers are flattened for the fully connected layer
\end{itemize}
\begin{center}
\small
\begin{tabular}{lc}
\hline
Layer & Output (channels, height, width) \\
\hline
input & (1, 28, 28) \\
conv1 + ReLU & (8, 26, 26) \\
max pool 2 x 2 & (8, 13, 13) \\
conv2 + ReLU & (16, 11, 11) \\
max pool 2 x 2 & (16, 5, 5) \\
flatten & 400 \\
fully connected & 10 \\
\hline
\end{tabular}
\end{center}
\end{frame}

%%%%%%%%%%%%%%%%%%%%%%%%%%%%%%%%%%%%%%%%%%%%%%%%%%%%%%%%%%
\begin{frame}[fragile]\frametitle{Counting Parameters}
\begin{itemize}
\item A filter is reused at every position of the image, so convolution layers need very few weights
\item Total: 5,258 parameters
\item A fully connected net (784 inputs, 128 hidden units, 10 outputs) needs 101,770: about 19 times more
\end{itemize}
\begin{center}
\small
\begin{tabular}{llr}
\hline
Parameter & Shape & Count \\
\hline
conv1 weights & (8, 1, 3, 3) & 72 \\
conv1 bias & (8) & 8 \\
conv2 weights & (16, 8, 3, 3) & 1,152 \\
conv2 bias & (16) & 16 \\
fc weights & (10, 400) & 4,000 \\
fc bias & (10) & 10 \\
\hline
total & & 5,258 \\
\hline
\end{tabular}
\end{center}
\end{frame}

%%%%%%%%%%%%%%%%%%%%%%%%%%%%%%%%%%%%%%%%%%%%%%%%%%%%%%%%%%
\begin{frame}[fragile]\frametitle{Training}
\begin{itemize}
\item Cross-entropy loss, the Adam optimizer, batches of 128 images, 3 passes (epochs) over the 60,000 images
\item Test accuracy is measured on the 10,000 images the network never trained on
\item It runs on a laptop CPU, and the seed is fixed so the numbers repeat on the same setup
\end{itemize}
\begin{lstlisting}
opt = torch.optim.Adam(model.parameters(), lr=1e-3)
for epoch in range(3):
    for i in range(0, len(xtr), 128):
        idx = perm[i:i + 128]
        opt.zero_grad()
        loss = F.cross_entropy(model(xtr[idx]), ytr[idx])
        loss.backward()
        opt.step()

# Test accuracy after epochs 1, 2, 3:
# 0.9518   0.9695   0.9770
\end{lstlisting}
\end{frame}

%%%%%%%%%%%%%%%%%%%%%%%%%%%%%%%%%%%%%%%%%%%%%%%%%%%%%%%%%%
\begin{frame}[fragile]\frametitle{CNN Against a Fully Connected Net}
\begin{itemize}
\item Same data, same optimizer, same 3 epochs
\item The CNN is wrong on 230 of the 10,000 test images; the fully connected net on 374
\item That is about 4 in 10 fewer mistakes, from a network with 19 times fewer parameters
\item One run with one seed: the gap is a good illustration, not a benchmark
\end{itemize}
\begin{center}
\small
\begin{tabular}{lrr}
\hline
 & Parameters & Test accuracy \\
\hline
Fully connected (784-128-10) & 101,770 & 96.26\% \\
Small CNN & 5,258 & 97.70\% \\
\hline
\end{tabular}
\end{center}
\end{frame}

%%%%%%%%%%%%%%%%%%%%%%%%%%%%%%%%%%%%%%%%%%%%%%%%%%%%%%%%%%
\begin{frame}[fragile]\frametitle{What the First Layer Learned}
\begin{itemize}
\item The 8 filters of the first layer after training, each a 3 x 3 grid of weights: red is negative, blue is positive
\item Nobody drew them: they started random and training shaped them
\item Some, such as f2 and f7, look like edge or line detectors: a strong row of one sign beside the opposite sign
\item Others are harder to read
\end{itemize}
\begin{center}
\includegraphics[width=0.95\linewidth]{dlcoep_cnn_filters.png}
\end{center}
\end{frame}

%%%%%%%%%%%%%%%%%%%%%%%%%%%%%%%%%%%%%%%%%%%%%%%%%%%%%%%%%%
\begin{frame}[fragile]\frametitle{Where It Goes Wrong}
\begin{itemize}
\item Easiest digits: 1 (99.4\%) and 0 (99.2\%); hardest: 7 (96.1\%) and 9 (96.2\%)
\item Most common confusions: 7 read as 2 (20 times), 4 as 9 (13), 9 as 3 (11), 2 as 8 (10)
\item The first 8 test mistakes, each labelled true digit and predicted digit
\item Some are hard even for a person, such as a 9 with an almost closed loop read as a 1
\end{itemize}
\begin{center}
\includegraphics[width=0.95\linewidth]{dlcoep_cnn_mistakes.png}
\end{center}
\end{frame}

%%%%%%%%%%%%%%%%%%%%%%%%%%%%%%%%%%%%%%%%%%%%%%%%%%%%%%%%%%
\begin{frame}[fragile]\frametitle{Quick Check: The Case Study}
\begin{block}{Think About It}
On MNIST the small CNN reached 97.70\% and the fully connected net 96.26\%, each in one run with one seed. Which conclusion is justified?
\begin{enumerate}
  \item[A.] CNNs always beat fully connected nets by about 1.4 points.
  \item[B.] In this run the CNN made fewer mistakes with far fewer parameters; a general claim needs more runs and more data.
  \item[C.] The fully connected net is overfitting.
  \item[D.] At 97.7\% the CNN is ready to read gauges in production.
\end{enumerate}
\end{block}
\end{frame}

%%%%%%%%%%%%%%%%%%%%%%%%%%%%%%%%%%%%%%%%%%%%%%%%%%%%%%%%%%
\begin{frame}[fragile]\frametitle{Quick Check: The Case Study}
\begin{block}{Answer}
\textbf{Correct: B.} One run supports only this case. A over-generalizes from one run, C is not shown by anything we measured (only test accuracy), and D needs a test on your own photographs, with your own lighting and fonts.
\end{block}
\end{frame}

%%%%%%%%%%%%%%%%%%%%%%%%%%%%%%%%%%%%%%%%%%%%%%%%%%%%%%%%%%
\begin{frame}[fragile]\frametitle{Try It Yourself}
\begin{itemize}
\item Script: \texttt{Code/\allowbreak dlcoep/\allowbreak session01\_cnn/\allowbreak cnn\_digits\_case\_study.py}, in the \texttt{dlcoep} conda environment
\item It downloads MNIST once (about 11 MB); with no network it falls back to scikit-learn's small 8 x 8 digits
\item Change one thing and predict the effect first: 4 filters in the first layer, 1 epoch, or a 5 x 5 filter
\item Which would you trust more on a real gauge: this 97.7\%, or a test on your own photographs?
\end{itemize}
\end{frame}
